% «Así lo preguntan en el examen» · Razonamiento lógico (libros 1 y 2). Sale sola antes del taller de cada lección.
% --- Libro 1 ---
\examende{pu-cm-1}{
\forma{Una operación combinada para calcular}{$-3^2 + 2(4 - 7)$. Potencia primero: $-9 + 2(-3) = -15$. Las opciones traen $3$ (quien hizo $(-3)^2$).}
\forma{Una historia de subidas y bajadas}{Temperaturas, pisos, saldos, goles a favor y en contra. «De $-6$ subió 15 y bajó 4»: $-6 + 15 - 4 = 5$.}
\forma{Distancias en la recta}{«¿Qué números están a 9 de $-2$?»: $-11$ y $7$. A veces dan una condición extra para quedarse con uno.}
\forma{¿Cuál expresión es siempre positiva o negativa?}{Con $a < 0 < b$: $ab$ es negativa; $a^2 b$ positiva; $a + b$ depende. Se resuelve con la regla de los signos, sin números.}}
\examende{pu-cm-2}{
\forma{La fracción de lo que queda}{«Gasta $\tfrac13$ y luego $\tfrac12$ del resto»: queda $\tfrac23 \cdot \tfrac12 = \tfrac13$. La opción trampa resta $\tfrac13 + \tfrac12$.}
\forma{Hallar el total desde una parte}{«$\tfrac35$ del tanque son 45 litros»: un quinto son 15, el tanque 75.}
\forma{Ordenar o comparar}{«¿Cuál es mayor, $\tfrac58$ o $\tfrac7{11}$?»: $55 < 56$, gana $\tfrac7{11}$.}
\forma{Un producto que se cancela en cadena}{$\tfrac12 \cdot \tfrac23 \cdot \tfrac34 \cdots \tfrac{9}{10} = \tfrac{1}{10}$.}}
\examende{pu-cm-3}{
\forma{La cifra que ocupa un lugar lejano}{«La cifra 50 de $0.\overline{27}$»: lugar par, es 7.}
\forma{Comparar números en formatos distintos}{$\tfrac38$, $0.4$ y $37\,\%$: todo a decimal ($0.375$, $0.40$, $0.37$).}
\forma{Porcentaje de una cantidad, al derecho y al revés}{«El 15\,\% de 80» (12) y «18 es el 15\,\% de…» (120).}
\forma{¿Qué porcentaje es?}{«12 de 48 estudiantes»: $\tfrac14 = 25\,\%$.}}
\examende{pu-cm-4}{
\forma{Simplificar con leyes de exponentes}{$\tfrac{2^{10} \cdot 4^3}{8^4} = 2^{10 + 6 - 12} = 16$.}
\forma{Contar cifras o ceros}{«¿Cuántas cifras tiene $2^6 \cdot 5^8$?»: $25 \cdot 10^6$, ocho cifras.}
\forma{Entre qué enteros está una raíz}{$\sqrt{200}$: entre 14 y 15, porque $196 < 200 < 225$.}
\forma{Qué pasa con el área si cambia el lado}{Lado $+10\,\%$: área $\times 1.21$, crece 21\,\%.}}
\examende{pu-ar-1}{
\forma{Varios gastos sucesivos y lo que sobra}{Se trabaja hacia atrás desde lo que quedó.}
\forma{«La mitad más…» repetido}{«Vende la mitad más 1, dos veces, y queda 3»: $(3 + 1)\cdot 2 = 8$; $(8 + 1)\cdot 2 = 18$.}
\forma{Rebotes, llenados y vaciados}{Cada paso multiplica por una fracción; al revés, se divide.}}
\examende{pu-ar-2}{
\forma{El dato que falta para un promedio}{«Promedio 4.0 en 5 notas y lleva 15.5»: necesita 4.5.}
\forma{Entra o sale alguien del grupo}{«El promedio de 9 es 20; llega uno y sube a 21»: el nuevo vale $210 - 180 = 30$.}
\forma{Dos grupos de distinto tamaño}{Se pesan: totales de cada grupo sobre el total de personas.}
\forma{Rapidez promedio de ida y vuelta}{No es el promedio de las rapideces: distancia total sobre tiempo total.}}
\examende{pu-ar-3}{
\forma{Una suma o resta con cifras tapadas}{Se resuelve columna por columna, de derecha a izquierda, anotando lo que se lleva.}
\forma{Letras que son cifras}{«$\overline{AB} + \overline{BA} = 132$»: $11(A + B) = 132$, $A + B = 12$.}
\forma{Invertir las cifras}{«Al invertirlas el número aumenta 36»: $9(b - a) = 36$.}}
\examende{pu-ar-4}{
\forma{«¿Cuántos sacar para estar seguros?»}{Peor caso más uno. Con 4 colores, para 2 iguales: 5.}
\forma{«Seguro de uno de un tipo específico»}{El peor caso agota primero todos los demás.}
\forma{El mínimo que comparte dos características}{«De 30, 20 usan gafas y 18 reloj»: al menos $38 - 30 = 8$ usan los dos.}
\forma{Máximo o mínimo con suma fija}{Para agrandar uno, se achican los demás.}}
\examende{pu-op-1}{
\forma{Calcular con una regla nueva}{$a \star b = 2a - b^2$; $3 \star (-2) = 6 - 4 = 2$.}
\forma{Dos operadores en la misma expresión}{Se respeta el paréntesis: primero lo de adentro.}
\forma{¿Es conmutativa?}{Se compara $a \star b$ con $b \star a$; un contraejemplo basta para decir que no.}}
\examende{pu-op-2}{
\forma{Una ecuación con el operador}{$x \star 2 = 10$ con $a \star b = 3a + b$: $3x + 2 = 10$.}
\forma{Una tabla de doble entrada}{Se lee fila y columna; preguntan el neutro, un inverso o una operación encadenada.}
\forma{Una regla que se llama a sí misma}{$f(1) = 2$, $f(n + 1) = f(n) + 4$: se calculan los primeros y se busca el patrón.}}
\examende{pu-pr-1}{
\forma{Reparto con razón y total}{«90 en razón $2:7$»: partes de 10, 20 y 70.}
\forma{Razón y diferencia}{«Razón $3:5$ y uno tiene 14 más»: dos partes son 14.}
\forma{Razones encadenadas}{$a:b = 2:3$ y $b:c = 6:5$: se iguala $b$ ($4:6:5$).}
\forma{Escalas de mapas y planos}{Longitudes por $k$, áreas por $k^2$.}}
\examende{pu-pr-2}{
\forma{Regla de tres directa}{Más kilos, más precio: se conserva el precio unitario.}
\forma{Regla de tres inversa}{Más obreros, menos días: se conserva el producto.}
\forma{Compuesta: obreros, horas, días, obra}{Se calcula el trabajo total y se divide.}
\forma{Cambia la condición a mitad de camino}{«A los 5 días se van 3 obreros»: se cuenta el trabajo hecho y el que falta.}}
\examende{pu-pr-3}{
\forma{Dos que se encuentran o se persiguen}{De frente se suman las velocidades; en persecución se restan.}
\forma{Dos que trabajan juntos}{Se suman las fracciones de trabajo por hora.}
\forma{Llave que llena y desagüe que vacía}{La rapidez del desagüe se resta.}
\forma{Rapidez promedio}{Distancia total entre tiempo total.}}
\examende{pu-pr-4}{
\forma{Las tres preguntas clásicas}{Parte, porcentaje o total: identifica cuál de los tres falta.}
\forma{Porcentaje de un porcentaje}{«El 20\,\% del 50\,\%» es el 10\,\%.}
\forma{Porcentaje de cambio}{$\tfrac{\text{final} - \text{inicial}}{\text{inicial}}$: la base es el valor de antes.}
\forma{Mezclas}{Se suman las cantidades de sustancia y se divide por el total.}}
\examende{pu-pr-5}{
\forma{Dos cambios seguidos}{Se multiplican los factores: $+10\,\%$ y $-10\,\%$ es $\times 0.99$.}
\forma{Hallar el precio de antes}{Se divide por el factor: sin IVA, $\div 1.19$.}
\forma{¿Cuánto hay que subir para recuperar?}{Tras bajar 20\,\%, hay que subir 25\,\%.}
\forma{Interés simple o compuesto}{Simple: $C r t$. Compuesto: $C(1 + r)^t$.}}
\examende{pu-ec-1}{
\forma{Un enunciado para traducir y resolver}{«El triple de un número menos 4 es 20»: $3x - 4 = 20$.}
\forma{Solo piden la ecuación}{Las opciones son cuatro ecuaciones: se elige la que traduce bien, sin resolver.}
\forma{Números consecutivos}{Se nombra el del medio para que la suma quede limpia.}
\forma{Perímetros y medidas}{«El largo es el doble del ancho y el perímetro 36».}}
\examende{pu-ec-2}{
\forma{Edades en dos momentos}{Tabla antes / hoy / después; la diferencia no cambia.}
\forma{Cabezas y patas, monedas, boletas}{Se supone que todos son del tipo barato y se reparte la diferencia.}
\forma{Mezclas de concentraciones o precios}{Se conserva la sustancia o el costo total.}}
\examende{pu-ec-3}{
\forma{Sistema que pide una combinación}{Si piden $x + y$, prueba sumar las ecuaciones antes de despejar.}
\forma{Dos compras con los mismos artículos}{«2 de A y 3 de B cuestan…»: reducción.}
\forma{Una desigualdad con contexto}{«¿Desde cuántos minutos conviene el plan?»: ojo con $>$ y $\ge$.}
\forma{Contar enteros en un intervalo}{$-3 \le x < 4$ tiene 7 enteros.}}
\examende{pu-nu-1}{
\forma{Contar divisores}{$72 = 2^3 \cdot 3^2$: $(3 + 1)(2 + 1) = 12$.}
\forma{Una cifra que falta para que sea divisible}{Criterios del 3, 9, 4, 11.}
\forma{¿Cuál es primo?}{Se prueban primos hasta la raíz.}
\forma{Hacer un cuadrado perfecto}{Se completan los exponentes impares.}}
\examende{pu-nu-2}{
\forma{Coincidencias en el tiempo}{Luces, buses, campanas: MCM.}
\forma{Partir en trozos iguales lo más grandes posible}{Varillas, baldosas, ramos: MCD.}
\forma{El menor número que deja cierto residuo}{«El menor número de dos cifras que deja residuo 2 al dividir entre 3, 4 y 5»: MCM $+ 2 = 62$.}
\forma{Producto, MCD y MCM}{$\text{MCD} \cdot \text{MCM} = a \cdot b$.}}
\examende{pu-nu-3}{
\forma{La última cifra de una potencia}{Ciclo de la base y residuo del exponente.}
\forma{El día de la semana dentro de $N$ días}{$N \bmod 7$.}
\forma{El residuo de una potencia}{Los residuos se repiten en ciclo.}
\forma{¿Es posible? (paridad)}{Suma de impares, repartos, tableros: mira si el total es par o impar.}}
\examende{pu-ge-1}{
\forma{Complementos y suplementos con una condición}{«El suplemento es el triple del complemento»: ecuación.}
\forma{Paralelas con una línea quebrada}{Se traza una paralela auxiliar por el quiebre.}
\forma{Ángulos de un triángulo o un polígono}{Suma $180^\circ$; polígono $(n - 2)180^\circ$; exteriores $360^\circ$.}
\forma{Las manecillas del reloj}{El minutero avanza $6^\circ$ por minuto; el horario, $0.5^\circ$.}}
\examende{pu-ge-2}{
\forma{Escaleras, cables, barcos}{Un triángulo rectángulo escondido en una historia; casi siempre una terna conocida.}
\forma{Diagonales}{De un rectángulo, un cuadrado ($\ell\sqrt2$) o una caja.}
\forma{Altura de un isósceles o equilátero}{La altura parte la base en dos.}
\forma{¿Se puede formar el triángulo? ¿De qué tipo es?}{Desigualdad triangular; comparar $c^2$ con $a^2 + b^2$.}}
\examende{pu-ge-3}{
\forma{Sombras a la misma hora}{Altura y sombra son proporcionales.}
\forma{Una paralela dentro de un triángulo}{Tales: el triángulo pequeño es semejante al grande.}
\forma{Áreas o volúmenes de figuras semejantes}{$k^2$ y $k^3$.}
\forma{Maquetas y mapas}{La escala es la razón de semejanza.}}
\examende{pu-ge-4}{
\forma{Figura compuesta}{Rectángulo más semicírculos, casa con techo triangular: se suma por piezas.}
\forma{Vueltas de una rueda}{Distancia entre la longitud de la circunferencia.}
\forma{Sector circular}{Fracción $\tfrac{\theta}{360}$ del círculo.}
\forma{Cambia una medida}{Radio doble: longitud doble, área cuádruple.}}
\examende{pu-ge-5}{
\forma{Región sombreada entre dos figuras}{Grande menos pequeña.}
\forma{Cuartos de círculo en las esquinas}{Juntos forman un círculo completo.}
\forma{Figura en una cuadrícula}{Encerrar en un rectángulo y restar triángulos.}
\forma{Triángulo con vértice móvil}{Misma base y misma altura: misma área.}}
\examende{pu-ge-6}{
\forma{Capacidad de un tanque o una caja}{Volumen y conversión a litros.}
\forma{Cubitos pintados}{3 caras: esquinas; 2: aristas; 1: centros; 0: interior.}
\forma{Cono frente a cilindro}{El cono es la tercera parte.}
\forma{Nivel que sube al sumergir un objeto}{Volumen $=$ área de la base por lo que sube.}}
% --- Libro 2 ---
\examende{pu-or-1}{
\forma{¿Quién ocupa cierto puesto?}{Se arma la fila con fijas, bloques y cadenas.}
\forma{¿Cuál afirmación es necesariamente cierta?}{Debe cumplirse en todas las filas posibles.}
\forma{Se agrega una pista en la pregunta}{«Si además Ana es tercera…»: se usa solo para esa pregunta.}}
\examende{pu-or-2}{
\forma{¿Quién está frente a…? ¿Quién a la derecha de…?}{Se fija a uno y se ubica el resto.}
\forma{¿Cuántas personas hay entre dos?}{Por un lado y por el otro suman el total menos dos.}
\forma{Conteo en mesas redondas}{$(n - 1)!$ maneras de sentar a $n$.}}
\examende{pu-or-3}{
\forma{Persona, profesión y ciudad}{Tabla con ✓ y ✗; las pistas cruzadas conectan categorías.}
\forma{Pistas con «mayor que», «no es»}{Las negativas tachan; las de orden encadenan.}
\forma{«No se puede determinar»}{Es una opción válida: si quedan dos soluciones, márcala.}}
\examende{pu-or-4}{
\forma{¿Cuál grupo es posible?}{Se prueba cada opción contra todas las condiciones.}
\forma{¿Quién tiene que estar?}{Se listan todos los grupos válidos y se mira qué se repite en todos.}
\forma{Deducciones con «si…, entonces…»}{Solo vale el contrarrecíproco.}}
\examende{pu-lo-1}{
\forma{Valor de verdad de una fórmula}{Se reemplazan V y F y se opera de adentro hacia afuera.}
\forma{Deducir los valores}{«$p \to q$ es falsa»: $p$ es V y $q$ es F.}
\forma{¿Cuál es tautología?}{Se busca si hay alguna forma de hacerla falsa.}}
\examende{pu-lo-2}{
\forma{Negar una frase con «todos» o «algunos»}{«Todos» $\leftrightarrow$ «algunos no»; «algunos» $\leftrightarrow$ «ninguno».}
\forma{Negar una frase con «y» u «o»}{De Morgan: se niegan las partes y se cambia el conector.}
\forma{Negar un condicional}{«Se cumple lo primero y no lo segundo».}}
\examende{pu-lo-3}{
\forma{¿Cuál frase es equivalente?}{El contrarrecíproco (o «no $p$ o $q$»).}
\forma{¿Qué se concluye?}{Ponens y tollens sí; afirmar el consecuente no.}
\forma{«Solo si», «necesario», «suficiente»}{Se traduce a una flecha antes de responder.}
\forma{Cadenas de condicionales}{$p \to q$ y $q \to r$ dan $p \to r$.}}
\examende{pu-lo-4}{
\forma{¿Qué conclusión se sigue?}{Diagramas: debe cumplirse en todo dibujo posible.}
\forma{¿Cuál conclusión NO se sigue?}{Basta un dibujo donde falle.}
\forma{Un caso particular}{«Todos los $A$ son $B$; Juan no es $B$»: Juan no es $A$.}}
\examende{pu-lo-5}{
\forma{¿Quién miente? ¿Quién lo hizo?}{Se buscan dos frases contradictorias o se supone y se sigue el hilo.}
\forma{«Solo uno dice la verdad»}{Se prueba cada persona como la veraz.}
\forma{La pregunta al guardián}{Acertijos de dos puertas: se razona qué respondería cada tipo.}}
\examende{pu-cj-1}{
\forma{Operar conjuntos dados por lista}{Unión, intersección, diferencia, complemento.}
\forma{¿Qué expresión describe la zona sombreada?}{Se nombra la zona con palabras («en $B$ y no en $A$») y se traduce.}
\forma{Contar subconjuntos}{$2^n$; con condiciones, se fijan elementos.}}
\examende{pu-cj-2}{
\forma{Dos conjuntos y «ninguno»}{$n(A \cup B) = $ total $-$ ninguno.}
\forma{Tres conjuntos}{Se llena del centro hacia afuera.}
\forma{¿Cuántos solo uno? ¿Cuántos exactamente dos?}{Se suman las zonas correspondientes del diagrama.}}
\examende{pu-co-1}{
\forma{Números con condiciones}{Cifras distintas, pares, que no empiecen en cero: primero la casilla restringida.}
\forma{Combinaciones de prendas o menús}{Producto de opciones.}
\forma{Claves y placas}{Con repetición: $n^k$.}}
\examende{pu-co-2}{
\forma{Ordenar personas u objetos}{$n!$; con algunos juntos, bloque.}
\forma{Anagramas}{Con letras repetidas se divide.}
\forma{Premios o cargos distintos}{$n(n - 1)(n - 2)\ldots$}}
\examende{pu-co-3}{
\forma{Comités y grupos}{$\binom nk$; con dos tipos de personas, se multiplica.}
\forma{Saludos, partidos, rectas, diagonales}{$\binom n2$.}
\forma{«Al menos uno»}{Total menos los que no tienen ninguno.}}
\examende{pu-co-4}{
\forma{Números con al menos una cifra dada}{Complemento.}
\forma{Múltiplos de uno u otro}{Inclusión-exclusión.}
\forma{Casos que no se solapan}{Se separa, se cuenta cada caso y se suma.}}
\examende{pu-co-5}{
\forma{Caminos en una cuadrícula}{$\binom{a + b}{a}$; con punto obligado, se multiplica por tramos.}
\forma{¿Cuántos triángulos (o cuadrados) hay en la figura?}{Se cuentan por tamaños.}
\forma{Caminos con una calle cerrada}{Se suman esquinas o se resta lo que pasa por la calle cerrada.}}
\examende{pu-pb-1}{
\forma{Dos dados}{Tabla de 36 casos.}
\forma{Urnas y bolsas}{Favorables sobre totales.}
\forma{Un número al azar con cierta propiedad}{Se cuentan los que la cumplen.}}
\examende{pu-pb-2}{
\forma{Dos extracciones sin devolver}{Se multiplican ramas; la segunda con lo que quedó.}
\forma{«Al menos una vez»}{$1 - P(\text{ninguna})$.}
\forma{Eventos independientes}{Producto de probabilidades.}}
\examende{pu-se-1}{
\forma{¿Qué número sigue?}{Diferencias, cocientes, mezclas, intercaladas.}
\forma{¿Qué número falta en el medio?}{La regla debe servir a los dos lados del hueco.}
\forma{El término de un lugar lejano}{$a_n = a_1 + (n - 1)d$.}
\forma{La suma de muchos términos}{$\tfrac{n(a_1 + a_n)}{2}$.}}
\examende{pu-se-2}{
\forma{¿Qué letra sigue?}{Letras a números.}
\forma{Pares o tríos de letras}{Cada posición sigue su propia regla.}
\forma{¿Qué figura sigue?}{Una característica a la vez: giro, cantidad, color.}}
\examende{pu-se-3}{
\forma{El elemento de un lugar lejano}{Residuo al dividir por la longitud del ciclo.}
\forma{¿Cuántas veces aparece?}{Ciclos completos más lo que sobra.}
\forma{Dos ciclos que coinciden}{MCM de las longitudes.}}
\examende{pu-cd-1}{
\forma{Codificar una palabra con la regla del ejemplo}{Se descubre el desplazamiento y se aplica.}
\forma{Decodificar}{Se aplica la regla al revés.}
\forma{Letras por números}{Posición en el alfabeto.}}
\examende{pu-cd-2}{
\forma{Un idioma inventado}{Se comparan frases y se aíslan palabras comunes.}
\forma{Números que se transforman}{Se busca la operación ($n^2 + 1$, doble menos uno…).}
\forma{Símbolos con significado oculto}{Igual que el idioma: lo que se repite se corresponde.}}
\examende{pu-es-1}{
\forma{¿Dónde termina?}{Tabla de dirección y posición.}
\forma{¿A qué distancia quedó del inicio?}{En cuadras: suma de diferencias; en línea recta: Pitágoras.}
\forma{¿Hacia dónde mira al final?}{Se cuentan giros a derecha e izquierda.}}
\examende{pu-es-2}{
\forma{¿Cuántos cubos hay?}{Por columnas o por capas.}
\forma{¿Cuántos faltan para completar un cubo?}{Total del cubo menos los que hay.}
\forma{¿Cuál es la vista desde arriba o de frente?}{Altura máxima por columna.}}
\examende{pu-es-3}{
\forma{¿Qué cara queda opuesta?}{Una cara en medio en línea recta.}
\forma{¿Cuál desarrollo arma el cubo?}{Se descartan los que tienen $2 \times 2$ o caras que se enciman.}
\forma{Dados}{Opuestas suman 7.}}
\examende{pu-es-4}{
\forma{Hoja doblada y perforada}{Se desdobla del último doblez al primero.}
\forma{Imagen en el espejo}{Izquierda y derecha se intercambian.}
\forma{El reloj en el espejo}{$12{:}00$ menos la hora que se ve.}}
\examende{pu-gr-1}{
\forma{Leer un valor o una diferencia}{Primero ejes y unidades.}
\forma{Porcentaje de aumento o disminución}{Sobre el valor inicial.}
\forma{De torta a cantidades}{Porcentaje por total; $1\,\% = 3.6^\circ$.}
\forma{¿En qué periodo cambió más?}{La barra o el tramo con mayor diferencia.}}
\examende{pu-gr-2}{
\forma{Calcular media, mediana o moda}{Mediana: ordenar primero.}
\forma{Tabla de frecuencias}{Media ponderada por las frecuencias.}
\forma{¿Qué pasa si cambia un dato?}{La media se mueve; la mediana, casi nunca.}}
\examende{pu-gr-3}{
\forma{¿Cuál afirmación es correcta según la gráfica?}{Se verifica cada una con un cálculo.}
\forma{¿Cuál NO se puede concluir?}{Causas, opiniones y generalizaciones.}
\forma{Dos gráficas a la vez}{Se cruza la información: porcentaje de una con total de la otra.}}

\InputIfFileExists{ejemplos-rl}{}{}
